% --- PAQUETES DE FORMATO Y LENGUAJE ---
\usepackage[utf8]{inputenc}
\usepackage[spanish, es-tabla]{babel} % Configuración para español
\usepackage[T1]{fontenc}
\usepackage{lmodern} % Fuente legible y profesional

% --- GEOMETRÍA DE PÁGINA (Ajustado a TFM estándar) ---
\usepackage[a4paper, top=3cm, bottom=3cm, left=3.5cm, right=2.5cm]{geometry}
\usepackage{fancyhdr} % Para encabezados y pies de página
\pagestyle{fancy}
\fancyhf{}
\fancyhead[L]{\small \leftmark} % Nombre del capítulo a la izquierda
\fancyhead[R]{\small \thepage}  % Número de página a la derecha
\renewcommand{\headrulewidth}{0.4pt} % Línea decorativa

% --- COLORES Y ENLACES ---
\usepackage[table]{xcolor}
\usepackage{hyperref}
\hypersetup{
    colorlinks=true,
    linkcolor=blue,
    filecolor=magenta,      
    urlcolor=cyan,
    citecolor=blue,
}

% --- GRÁFICOS Y TABLAS ---
\usepackage{graphicx}
\usepackage{booktabs} % Para tablas estilo profesional (sin líneas verticales)
\usepackage{caption}
\captionsetup{font=small, labelfont=bf} % Etiquetas de figuras en negrita

% --- CONFIGURACIÓN DE SECCIONES (Estilo Romano como tu PDF) ---
% Esto hace que los capítulos se vean como "SECCIÓN I", "SECCIÓN II", etc.
\usepackage{titlesec}
\titleformat{\chapter}[display]
  {\normalfont\huge\bfseries\centering}
  {\chaptertitlename\ \thechapter}
  {20pt}
  {\Huge}
\renewcommand{\thechapter}{\Roman{chapter}} % Cambia números 1, 2... por I, II...

% --- CONFIGURACIÓN DE CÓDIGO (VHDL para VICON) ---
\usepackage{listings}
\definecolor{codegreen}{rgb}{0,0.6,0}
\definecolor{codepurple}{rgb}{0.58,0,0.82}

\lstdefinelanguage{VHDL}{
   morekeywords={library,use,all,entity,is,port,in,out,end,architecture,of,begin,process,if,then,else,elsif,case,when,signal,variable,type,range,array,component,generic,map},
   morecomment=[l]--
}

\lstset{
    language=VHDL,
    basicstyle=\ttfamily\small,
    keywordstyle=\color{blue}\bfseries,
    commentstyle=\color{codegreen},
    stringstyle=\color{codepurple},
    numbers=left,
    numberstyle=\tiny\color{gray},
    stepnumber=1,
    numbersep=10pt,
    backgroundcolor=\color{gray!5},
    showspaces=false,
    showstringspaces=false,
    frame=single,           % Recuadro alrededor del código
    rulecolor=\color{black},
    breaklines=true,        % Corta líneas largas automáticamente
    captionpos=b            % Título del código abajo
}